\chapter[Analytic Theory]{Analytic Theory of the Universal Optimality Defect}\label{ch:sec:analytic-uod}

\section{Introduction}\label{ch:sec:introduction-analytic_uod}

This chapter develops three complementary perspectives on the Universal
Optimality Defect: algebraic (Section~\ref{ch:sec:algebra-uod}), differential (Section~\ref{ch:sec:differential-uod}), and
functional-analytic (Section~\ref{ch:sec:functional-uod}).  Each reveals different structural
properties that are used in the subsequent chapters.

\paragraph{Algebra of the UOD.}
The UOD inherits an algebraic structure from the logarithmic embedding.
\label{ch:sec:algebra-uod}


\chapter[Algebra of the UOD]{Algebra of the Universal Optimality Defect}\label{ch:sec:algebra-of-the-universal-optimality-defect}
\section{Introduction}\label{ch:sec:introduction-analytic_uod-2}
Part~\ref{ch:part:posterior} constructed the posterior manifold together with the
canonical logarithmic coordinate map
$F = Q\circ L, \qquad L(P)=\log P,$
where (Q) denotes the quotient projection removing the additive
logarithmic gauge symmetry. The Universal Optimality Defect (UOD) is the
intrinsic squared distance induced by this geometry:
$\mathcal{D}(P,Q)=\|F(P)-F(Q)\|^{2}.$
The purpose of this chapter is to derive an exact closed formula for
\(\mathcal{D}\) and establish its fundamental algebraic properties.
All later comparison theorems will be derived from this identity.
\section{Canonical Hyperplane Representation}\label{ch:sec:canonical-hyperplane-representation}
Part~\ref{ch:part:posterior} defines (Q) abstractly as the quotient map
\[ Q:\mathbb{R}^{n}\rightarrow
\mathbb{R}^{n}/\operatorname{span}{\mathbf1}. \]
For explicit computations we identify this quotient with the Euclidean
hyperplane
\[ H=\Bigl\{x\in\mathbb{R}^{n}:\sum_{i}
x_i=0\Bigr\}. \]
\begin{lemma}
\label{ch:lem:21-1}
The quotient space
\(\mathbb{R}^{n}/\operatorname{span}{\mathbf1}\) is
canonically isometric to (H). Under this identification,
\[ Q(x)=x-\bar x,\mathbf1, \qquad
\bar x=\frac1n\sum_{i=1}^{n} x_i. \]
\begin{proof}
Every equivalence class
\(x+\operatorname{span}{\mathbf1}\) contains a unique
vector with zero arithmetic mean, namely
\(x-\bar x,\mathbf1\). Uniqueness follows because adding a
constant changes the mean by that constant. Hence every quotient class
has a unique representative in (H), defining a linear isomorphism. The
Euclidean norm on (H) therefore induces the canonical quotient norm used
throughout Part~\ref{ch:part:posterior}.
\end{proof}
\end{lemma}
\section{Closed Formula}\label{ch:sec:closed-formula}
\begin{theorem}
\label{ch:thm:21-2}
Let
\[ a_i=\log\frac{P_i}{Q_i}, \qquad
\bar a=\frac1n\sum_{i} a_i. \]
Then
\[ \boxed{
\mathcal{D}(P,Q)
=
\sum_{i=1}^{n}(a_i-\bar a)^2.
} \]
Equivalently,
\[ \boxed{
\mathcal{D}(P,Q)
=
\sum_i a_i^2
-
\frac1n\left(\sum_i a_i\right)^2.
} \]
\begin{proof}
By definition,
$\mathcal{D}(P,Q) = \|Q(\log P-\log Q)\|^{2}.$
Set \(a=\log P-\log Q\). By Lemma~\ref{ch:lem:21-1},
$Q(a)=a-\bar a,\mathbf1.$
Taking the Euclidean norm gives
\[ \mathcal{D}(P,Q) = \|a-\bar a,\mathbf1\|^{2},
\]
whose expansion is
\[ \sum_{i}(a_i-\bar a)^{2} = \sum_{i}
a_i^{2}-\frac1n\left(\sum_{i}
a_i\right)^{2}. \]
\end{proof}
\end{theorem}
\section{Equivalent Forms}\label{ch:sec:equivalent-forms}
\begin{corollary}
\label{ch:cor:21-3}
The Universal Optimality Defect admits the variance representation
\[ \boxed{
\mathcal{D}(P,Q)
=
n\,
\operatorname{Var}_{U_n}
\!\left(
\log\frac{P_i}{Q_i}
\right),
} \]
where \(\operatorname{Var}_{U_n}\) denotes variance with respect
to the uniform distribution on the coordinate index set.
\end{corollary}
\section{Fundamental Properties}\label{ch:sec:fundamental-properties}
\begin{proposition}
\label{ch:prop:21-4}
For all posterior distributions (P,Q):
\begin{itemize}
  \item \(\mathcal{D}(P,Q)\ge0\);
  \item \(\mathcal{D}(P,Q)=0\iff P=Q\);
  \item \(\mathcal{D}(P,Q)=\mathcal{D}(Q,P)\);
  \item \(\mathcal{D}\) is invariant under addition of a common
\end{itemize}
    constant in logarithmic coordinates.
\begin{proof}
Items (1) and (3) follow from the squared Euclidean norm.
If \(\mathcal{D}(P,Q)=0\), then all centered logarithmic ratios
vanish, so
$\log\frac{P_i}{Q_i}=c$
for every (i). Hence \(P_i=e^{c}Q_i\). Since both distributions sum to one,
$1=\sum_{i}P_i=e^{c}\sum_{i}Q_i=e^{c},$
thus (c=0), proving (P=Q).
Item (4) follows because (Q) annihilates constant vectors.
\end{proof}
\end{proposition}
\section{Interpretation}\label{ch:sec:interpretation}
The UOD is the intrinsic quadratic energy of the logarithmic
probability-ratio vector after quotienting the additive gauge symmetry.
Unlike Kullback--Leibler divergence, which averages logarithmic ratios
with respect to a probability measure, the UOD measures their Euclidean
dispersion in the canonical quotient geometry determined by Part~\ref{ch:part:posterior}.
\section{Outlook}\label{ch:sec:outlook}
Part~\ref{ch:part:variational} develops the mathematical theory of the UOD:
\begin{itemize}
  \item differential structure;\label{ch:sec:differential-uod}
  \item local Riemannian approximation;
  \item functional-analytic properties;\label{ch:sec:functional-uod}
  \item global comparison theorems with Total Variation, Hellinger, $\chi^2$ and
\end{itemize}
    Kullback--Leibler divergence;
\begin{itemize}
  \item applications to entropy and security functionals.
\end{itemize}


\paragraph{Running example: geo-indistinguishability.}
The variance representation $\mathcal D(P,Q)=n\operatorname{Var}_{U_n}(a)$ (Corollary~\ref{ch:cor:21-3}) is particularly illuminating for geo-indistinguishability (Section~\ref{ch:sec:geo-indistinguishability-and-location-privacy}): the log-likelihood ratio $a(x,y)=\log(P(x)/Q(y))$ becomes the Euclidean distance between two points on the plane, and the UOD is proportional to the squared spatial distance---establishing a direct link between geometric privacy and Euclidean geometry.

\chapter[Differential Theory of the UOD]{Differential Theory of the Universal Optimality Defect}\label{ch:sec:differential-theory-of-the-universal-optimality-defect}

\section{Introduction}\label{ch:sec:introduction-analytic_uod-3}

The preceding chapters established an exact algebraic formula for the Universal
Optimality Defect (UOD). The present chapter develops its differential
theory.

Throughout, let

\[
F=Q\circ L:\mathcal P\rightarrow F(\mathcal P)
\]

be the global logarithmic embedding of Part~\ref{ch:part:posterior} and

\[
\mathcal D(P,Q)=\|F(P)-F(Q)\|^{2}.
\]

Since \(F\) is a global diffeomorphism and the pullback metric is defined
by \(g=F^{*}g_E\), all differential properties of the UOD can be derived
directly from Euclidean calculus.

\section{Differential of the Logarithmic Embedding}\label{ch:sec:differential-of-the-logarithmic-embedding}

For \(P\in\mathcal P\),

\[
DF_P=Q\circ DL_P.
\]

By Part~\ref{ch:part:posterior}, \(DF_P\) is a linear isomorphism on every tangent space.
Consequently,

\[
g_P(u,v)=\langle DF_Pu,DF_Pv\rangle.
\]

\section{First Variation}\label{ch:sec:first-variation}

\begin{theorem}[First Variation of the UOD]
\label{ch:thm:22-1}
For every smooth curve \(P(t)\),

\[
\frac{d}{dt}\mathcal D(P(t),Q)=
2\langle DF_{P(t)}P'(t),F(P(t))-F(Q)\rangle.
\]
\end{theorem}

\begin{proof}
Differentiate the Euclidean squared norm and apply the chain rule.
$\blacksquare$
\end{proof}

\section{Characterization of Critical Points}\label{ch:sec:characterization-of-critical-points}

\begin{theorem}[Unique Critical Point]
\label{ch:thm:22-2}
The only critical point of

\[
P\mapsto\mathcal D(P,Q)
\]

is \(P=Q\).
\end{theorem}

\begin{proof}
If \(P\) is critical then \(d\mathcal D_P(v)=0\) for every tangent
vector \(v\). By Theorem~\ref{ch:thm:22-1},

\[
\langle DF_Pv,F(P)-F(Q)\rangle=0
\]

for every \(v\). Since \(DF_P\) is an isomorphism, its image is the
whole tangent space of \(F(\mathcal P)\). Hence \(F(P)-F(Q)=0\).
The injectivity of \(F\) (Part~\ref{ch:part:posterior}) implies \(P=Q\). $\blacksquare$
\end{proof}

\section{Gradient}\label{ch:sec:gradient}

\begin{corollary}[Riemannian Gradient]
\label{ch:cor:22-3}
The Riemannian gradient is uniquely determined by

\[
g_P(\nabla_g\mathcal D,v)=2\langle DF_Pv,F(P)-F(Q)\rangle.
\]

Equivalently,

\[
\nabla_g\mathcal D = 2\,DF_P^{-1}(F(P)-F(Q)).
\]
\end{corollary}

\begin{proof}
Using the pullback metric \(g_P(u,v)=\langle DF_Pu,DF_Pv\rangle\), the
defining identity of the Riemannian gradient gives

\[
DF_P(\nabla_g\mathcal D)=2(F(P)-F(Q)).
\]

Invertibility of \(DF_P\) completes the proof. $\blacksquare$
\end{proof}

\section{Hessian at the Reference Posterior}\label{ch:sec:hessian-at-the-reference-posterior}

\begin{theorem}[Hessian at the Optimum]
\label{ch:thm:22-4}
At the reference point \(P=Q\),

\[
\operatorname{Hess}_g(\mathcal D)_P = 2I.
\]
\end{theorem}

\begin{proof}
Near \(P\),

\[
F(P+\delta) = F(P)+DF_P\delta+o(\|\delta\|).
\]

Therefore,

\[
\mathcal D(P+\delta,P) = \|DF_P\delta\|^{2}+o(\|\delta\|^{2}).
\]

Using the definition of the pullback metric,
\(\|DF_P\delta\|^{2}=g_P(\delta,\delta)\), so

\[
\mathcal D(P+\delta,P) = g_P(\delta,\delta)+o(\|\delta\|^{2}).
\]

Hence the quadratic term is represented by the identity bilinear form,
whose Hessian is \(2I\). $\blacksquare$
\end{proof}

\section{Strict Convexity Near the Optimum}\label{ch:sec:strict-convexity-near-the-optimum}

\begin{corollary}[Local Convexity]
\label{ch:cor:22-5}
The reference posterior is a strict local minimizer of the UOD.
\end{corollary}

\begin{proof}
The Hessian at \(P\) equals \(2I\), which is positive definite.
$\blacksquare$
\end{proof}

\section{Local Quadratic Model}\label{ch:sec:local-quadratic-model}

Combining the previous results,

\[
\boxed{
\mathcal D(P+\delta,P)
=
\|\delta\|_g^2
+
o(\|\delta\|_g^2).
}
\]

Thus the UOD is locally identical to the squared Riemannian distance
induced by the pullback metric.

\section{Consequences}\label{ch:sec:consequences}

The differential theory establishes:
\begin{itemize}
\item smoothness of the UOD;
\item uniqueness of its critical point;
\item explicit gradient;
\item exact Hessian at the optimum;
\item strict local convexity;
\item local quadratic approximation.
\end{itemize}

These properties are intrinsic consequences of the logarithmic embedding
and require no additional probabilistic assumptions.

\section{Outlook}\label{ch:sec:outlook-analytic_uod}

The next chapter compares this local quadratic structure with the
classical Fisher information metric, identifying the precise
relationship between the WIS geometry and standard information geometry.


\chapter[Functional-Analytic Properties]{Functional-Analytic Properties of the Universal Optimality Defect}\label{ch:sec:functional-analytic-properties-of-the-universal-optimal}
\section{Introduction}\label{ch:sec:introduction-analytic_uod-4}
This chapter studies the global analytical properties of the Universal
Optimality Defect (UOD). Throughout,
\[
F=Q\circ L:\mathcal P^\circ\rightarrow F(\mathcal P^\circ),
\qquad \mathcal D(P,Q)=\|F(P)-F(Q)\|^2,
\]
where \(\mathcal P^\circ\) denotes the interior of the
probability simplex. Part~\ref{ch:part:posterior} proves that \(F\) is a global
diffeomorphism onto its image. The results below are consequences of
\section{Smoothness}\label{ch:sec:smoothness}
\begin{theorem}[Smoothness]
\label{ch:thm:23-1}
\[
\mathcal D:\mathcal P^\circ\times\mathcal P^\circ\rightarrow\mathbb R
\]
is \(C^\infty\).
\end{theorem}
\begin{proof} The logarithm is smooth on \((0,\infty)\), \(Q\) is
linear, and the squared Euclidean norm is polynomial. Therefore
\(\mathcal D\) is a composition of smooth maps. 
\end{proof}
%% Explanatory comment: Smoothness of the UOD implies that the Universal Optimality Defect is infinitely differentiable, which allows the use of differential calculus for its analysis. Consequently, gradient-based optimization methods and Taylor expansions are applicable.
\section{Positivity and Separation}\label{ch:sec:positivity-and-separation}
\begin{theorem}[Positivity]
\label{ch:thm:23-2}
$\mathcal D(P,Q)\ge0,$
with equality iff \(P=Q\).
\end{theorem}
\begin{proof} Non-negativity follows from the Euclidean norm. If
\(\mathcal D(P,Q)=0\), then \(F(P)=F(Q)\). By the results of Part~\ref{ch:part:posterior}, \(F\) is
injective, hence \(P=Q\). The converse is immediate. 
\end{proof}
%% Explanatory comment: Positivity ensures that the UOD is a genuine measure of discrepancy between probability distributions, vanishing only when the distributions coincide. This property is essential for the UOD to function as a valid distance-like functional.
\section{Continuity}\label{ch:sec:continuity}
\begin{theorem}[Continuity]
\label{ch:thm:23-3}
If \(P_n\to P\) and \(Q_n\to Q\), then
$\mathcal D(P_n,Q_n)\to\mathcal D(P,Q).$
\end{theorem}
\begin{proof} The UOD is the composition of the smooth embedding
\(F\) with the squared Euclidean norm, both of which are \(C^\infty\)
on the interior of the simplex.  Hence \(\mathcal D\) is smooth.
\end{proof}
%% Explanatory comment: Joint continuity guarantees that the UOD is stable under simultaneous small perturbations of both arguments. In practice this means that nearby probability distributions produce nearby UOD values, a prerequisite for any reliable numerical computation.
\section{Local Lipschitz Continuity}\label{ch:sec:local-lipschitz-continuity}
\begin{theorem}[Local Lipschitz]
\label{ch:thm:23-4}
For every compact set \(K\subset\mathcal P^\circ\), there exists
\(L_K>0\) such that
\[ |\mathcal D(P_1,Q_1)-\mathcal D(P_2,Q_2)|
\le L_K(\|P_1-P_2\|+\|Q_1-Q_2\|) \]
for all \(P_i,Q_i\in K\).
\end{theorem}
\begin{proof} Since \(\mathcal D\) is \(C^{1}\), its differential is
bounded on \(K\times K\). The multivariable mean-value theorem
gives the estimate. 
\end{proof}
%% Explanatory comment: Local Lipschitz continuity provides a quantitative bound on how much the UOD can change under perturbations: small changes in the input probabilities produce proportionally small changes in the defect value. This estimate is crucial for convergence analysis of iterative algorithms.
\section{An Auxiliary Linear Lemma}\label{ch:sec:an-auxiliary-linear-lemma}
The proof of properness relies on a purely linear fact.
\begin{lemma}[Growth in the Quotient Space]
\label{ch:lem:23-5}
Let
\[ Q(a)=a-\bar a\mathbf 1, \qquad
\bar a=\frac1n\sum_i a_i. \]
Suppose a sequence \(a^{(k)}\in\mathbb R^n\) satisfies
\[ \operatorname{diam}(a^{(k)}) = \max_i
a_i^{(k)}-\min_i a_i^{(k)}
\longrightarrow\infty. \]
Then
$\|Q(a^{(k)})\|\longrightarrow\infty.$
\end{lemma}
\begin{proof} The vector \(Q(a)\) is obtained by subtracting the mean from
every coordinate. Subtracting a constant does not change pairwise
differences:
$(Q(a))_i-(Q(a))_j=a_i-a_j.$
Hence
\[ \operatorname{diam}(Q(a)) = \operatorname{diam}(a). \]
For every zero-mean vector \(x\),
\[ \|x\| \ge \frac{\operatorname{diam}(x)}{\sqrt 2}, \]
because two coordinates differ by at most \(\sqrt 2\|x\|\).
Therefore,
\[ \|Q(a^{(k)})\| \ge
\frac{\operatorname{diam}(a^{(k)})}{\sqrt 2}
\longrightarrow\infty. \]
\end{proof}
%% Explanatory comment: This lemma establishes that when the spread of logarithmic coordinates diverges, the norm in the quotient space diverges as well. It is the key linear ingredient for proving that the UOD diverges near the boundary of the simplex.
\section{Properness}\label{ch:sec:properness}
\begin{theorem}[Properness]
\label{ch:thm:23-6}
Fix \(Q\in\mathcal P^\circ\). If \(P_n\) approaches
the boundary of the simplex, then
$\mathcal D(P_n,Q)\to\infty.$
\end{theorem}
\begin{proof} Approaching the boundary implies that at least one coordinate
satisfies
$P_{n,i}\to0,$
hence
$\log\frac{P_{n,i}}{Q_i}\to-\infty.$
Since \(\sum_i P_{n,i}=1\), at least one coordinate remains bounded
away from zero along a subsequence. Therefore the diameter of the
logarithmic ratio vector
$a^{(n)}=\log P_n-\log Q$
tends to infinity. Lemma~\ref{ch:lem:23-5} yields
$\|Q(a^{(n)})\|\to\infty,$ which is precisely \(\mathcal D(P_n,Q)\to\infty\). 
\end{proof}
%% Explanatory comment: Properness prevents minimizing sequences from escaping to the boundary of the simplex: if a sequence approaches the boundary, the UOD diverges to infinity. This guarantees that every minimization problem formulated with the UOD admits a solution in the interior.
\section{Compact Sublevel Sets}\label{ch:sec:compact-sublevel-sets}
\begin{corollary}[Compact Sublevel Sets]
\label{ch:cor:23-7}
$S_c=\{P\in\mathcal P^\circ:\mathcal D(P,Q)\le c\}$
is compact in the relative topology of the simplex.
\end{corollary}
\begin{proof} Continuity implies that \(S_c\) is closed. Properness prevents
sequences in \(S_c\) from approaching the boundary. Since the simplex is
compact, \(S_c\) is compact. 
\end{proof}
%% Explanatory comment: Compact sublevel sets mean that every set of distributions whose UOD from a fixed reference is bounded below a threshold is necessarily compact. Consequently, any continuous objective defined on such sublevel sets attains its minimum, a property essential for the existence of optimal posterior distributions.
\section{Topological Equivalence}\label{ch:sec:topological-equivalence}
\begin{theorem}[Topological Equivalence]
\label{ch:thm:23-8}
The topology induced by the UOD coincides with the manifold topology on
\(\mathcal P^\circ\).
\end{theorem}
\begin{proof} Part~\ref{ch:part:posterior} proves that \(F\) is a homeomorphism onto its image.
Since the UOD is the squared Euclidean distance in that image,
convergence in UOD is equivalent to Euclidean convergence of \(F(P)\),
hence to manifold convergence. 
\end{proof}
%% Explanatory comment: Topological equivalence confirms that the UOD induces the same notion of convergence as the underlying manifold topology. Thus the UOD is not merely a convenient computational tool but a geometrically faithful representation of the proximity structure of the probability simplex.
\section{Summary}\label{ch:sec:summary}
The Universal Optimality Defect is:
\begin{itemize}
\item infinitely differentiable;
\item positive definite;
\item jointly continuous;
\item locally Lipschitz on compact subsets;
\item proper on the interior simplex;
\item endowed with compact sublevel sets;
\item topologically equivalent to the intrinsic WIS geometry.
\end{itemize}
These properties establish the UOD as a globally well-behaved analytical
functional, providing the foundation for the comparison theorems
developed in Part~\ref{ch:part:variational}.
